\slajdid{09_2-s037}
% element: 09_2-s037:e01
% element: 09_2-s037:e02
% element: 09_2-s037:e03
% element: 09_2-s037:e04
\begin{frame}[label=09_2-s037]{CMOS: pet oblasti karakteristike}
\begin{columns}[T,onlytextwidth]\column{.55\textwidth}\centering {\RazaviSetup
\begin{circuitikz}[x=1cm,y=1cm,font=\fontsize{10}{12}\selectfont]
\node[rz nmos] (N) at (.28,.65) {};
\node[rz pmos] (P) at (.28,2.2) {};
\coordinate (O) at (N.D |- 0,1.43);
\draw[wire] (N.D)--(O)--(P.D);
\draw[wire] (N.S)--++(0,-.20) coordinate (GND);
\RazaviGround{GND}
\draw[wire] (P.S)--++(0,.22) coordinate (VDD);
\RazaviSupply{VDD}{V_{DD}}
\draw[wire] (O)--++(.60,0) node[rz terminal]{} node[right=3pt]{$V_O$};
\fill (O) circle(1.15pt);
\node[right=3pt] at (N.east) {$\mathrm{N}$};
\node[right=3pt] at (P.east) {$\mathrm{P}$};
\coordinate (I) at (-.7,1.43);
\draw[wire] (N.G)--(I |- N.G)--(I |- P.G)--(P.G);
\draw[wire] (I)--++(-.25,0) node[rz terminal]{} node[left=3pt]{$V_I$};
\fill (I) circle(1.15pt);
\end{circuitikz}}
\hspace{5mm}\centering \begin{tikzpicture}[x=.58cm,y=.52cm,font=\sitneoznake,>=Latex]
\draw[->](-.15,0)--(4.3,0)node[right]{$V_I$};\draw[->](0,-.15)--(0,3.7)node[above]{$V_O$};
\draw(0,3.1)node[left]{$V_{OH}$}--(.75,3.1)..controls(1.2,3.1)and(1.7,2.8)..(1.7,2.3)--(1.7,1.05)..controls(1.7,.48)and(2.2,0)..(2.9,0)--(3.75,0);
\node[left]at(0,0){$V_{OL}$};\node[below]at(3.85,0){$V_{DD}$};\draw[dashed](1.37,2.92)--(1.37,0)node[below left]{$V_{IL}$};\draw[dashed](2.16,.38)--(2.16,0)node[below]{$V_{IH}$};
\draw(1.1,3.17)--(1.7,2.57);\node[right]at(1.43,3.1){$a=-1$};\draw(1.94,.61)--(2.51,.04);\node[right]at(2.28,.85){$a=-1$};\draw[dashed](0,0)--(1.95,1.95)node[right]{$V_S$};
\end{tikzpicture}
\column{.41\textwidth}\begin{enumerate}
\item N zakočen; P omski.
\item N zasićen; P omski.
\item N zasićen; P zasićen.
\item N omski; P zasićen.
\item N omski; P zakočen.
\end{enumerate}\end{columns}\[V_{O(IH)}=\frac{V_{DD}-2V_{Tn}}8\le V_{IL}=\frac{3V_{DD}+2V_{Tn}}8\]
\[V_{IH}=\frac{5V_{DD}-2V_{Tn}}8\le V_{O(IL)}=\frac{7V_{DD}+2V_{Tn}}8\]
\centering Oba uslova daju $V_{DD}\geq-2V_{Tn}$.
\end{frame}
